\documentclass[twocolumn,a4paper,10pt]{article}

\usepackage{amsmath,amssymb}
\usepackage{booktabs}
\usepackage{hyperref}
\usepackage{graphicx}
\usepackage{url}
\usepackage[numbers,sort&compress]{natbib}
\usepackage{multirow}
\usepackage{array}
\usepackage{xcolor}
\usepackage{geometry}
\usepackage{caption}
\usepackage{subcaption}
\usepackage{microtype}
\geometry{a4paper, top=2cm, bottom=2cm, left=1.8cm, right=1.8cm, columnsep=0.6cm}

\setlength{\parskip}{2pt}

\title{\textbf{PEARL: Parameter-Efficient Adaptation with
Retrieval-Augmented Learning for Molecular Property Prediction}\\[4pt]
{\large Integrating LoRA, 3D Molecular Representations, and Chemical
Neighbourhood Retrieval across ChemBERTa, MolFormer, and Uni-Mol}}

\author{
Kavya Agrawal$^{1}$ \quad Harsha Harod$^{1}$ \quad Ines Simeone$^{2}$ \quad
Michele Ceccarelli$^{3}$ \quad Halima Bensmail$^{5}$ \quad
Sukrit Gupta$^{1}$ \quad Raghvendra Mall$^{4,5}$\\[4pt]
{\small $^{1}$Center for Applied Research in Data Science, IIT Ropar, Punjab, India}\quad
{\small $^{2}$Humanitas Research, Rozzano, Italy}\\
{\small $^{3}$Sylvester Cancer Center, University of Miami, Florida, U.S.A}\\
{\small $^{4}$Technology Innovation Institute, Abu Dhabi, U.A.E.}\quad
{\small $^{5}$Qatar Computing Research Institute, HBKU, Doha, Qatar}\\[2pt]
\texttt{ramall@hbku.edu.qa, raghvendramall@ieee.org}
}

\date{Received: 10 March 2026 \quad|\quad Revised: \quad|\quad Accepted:}

\begin{document}

\maketitle

\noindent\textbf{Keywords:} PEARL $|$ Parameter-Efficient Finetuning $|$ Chemical Language Models $|$
Low-Rank Adaptation $|$ 3D Molecular Representations $|$ Retrieval-Augmented Generation $|$
Molecular Properties $|$ Computational Drug Discovery

\begin{abstract}
\noindent\textbf{Motivation:}
Chemical language models (CLMs) such as ChemBERTa and MolFormer enable compound property
prediction from SMILES sequences, but their computational demands and reliance on 1D string
representations limit performance on tasks that depend on three-dimensional molecular geometry.
We introduce PEARL (\textbf{P}arameter-\textbf{E}fficient \textbf{A}daptation with
\textbf{R}etrieval-augmented \textbf{L}earning), a framework that (i)~integrates Low-Rank
Adapters (LoRA) with three distinct model families---ChemBERTa, MolFormer, and Uni-Mol---and
(ii)~augments predictions with a Retrieval-Augmented Feature Enhancement (RAFE) pipeline that
retrieves chemical neighbourhood context from a 249,455-molecule ZINC-250k knowledge base.

\noindent\textbf{Results:}
Uni-Mol, which encodes explicit 3D atomic coordinates via a SE(3)-equivariant Transformer,
achieves the highest performance on the BACE inhibitor task (MCC~$= 0.623$, AUC~$= 0.891$),
representing an 8\% MCC improvement over MolFormer.
The RAFE framework establishes a new state of the art on ClinTox toxicity prediction
(CT\_TOX MCC~$= 0.913$), surpassing the previous best of $0.875$.
On blood-brain barrier permeability prediction (BBBP), RAFE with MolFormer achieves
MCC~$= 0.798$, matching end-to-end LoRA (MCC~$= 0.797$) while providing an interpretable
chemical-space context signal.
All models achieve 75--96\% parameter reduction via LoRA.

\noindent\textbf{Availability:}
Code is available at \href{https://github.com/raghvendra5688/PEARL}{github.com/raghvendra5688/PEARL}.
\end{abstract}

%% ─────────────────────────────────────────────────────────────────────────────
\section{Introduction}
%% ─────────────────────────────────────────────────────────────────────────────

The study of molecular structure--property relationships has been central to advances in
computational drug discovery. Although traditional quantitative structure--activity relationship
(QSAR) methods rely on handcrafted molecular descriptors, deep learning techniques, particularly
transformer-based architectures~\cite{vaswani2017attention}, enable end-to-end learning of
molecular representations directly from their Simplified Molecular Input Line Entry System
(SMILES) sequences. This has catalysed the development of chemical language models (CLMs)
pretrained on massive molecular corpora, now rivalling experimental approaches in predicting
properties such as toxicity, solubility, and bio-activity~\cite{ahmad2022chemberta2,huang2020deeppurpose}.

Despite their promise, two critical challenges persist: (1)~CLMs process SMILES as flat token
sequences and therefore lack explicit inductive biases for three-dimensional molecular geometry,
which is fundamental to binding affinity and enzyme inhibition; and (2)~the computational cost
of fine-tuning $\mathcal{O}(10^8)$-parameter CLMs for each specialised task remains prohibitive
in resource-constrained settings.

This work addresses both challenges through a systematic study of three model families:
MolFormer~\cite{ross2022molformer}, ChemBERTa-2~\cite{ahmad2022chemberta2}, and
Uni-Mol~\cite{zhou2023unimol}---a SE(3)-equivariant Transformer that operates directly on
explicit 3D atomic coordinates.
We apply Low-Rank Adaptation (LoRA)~\cite{hu2021lora} to all three families, dramatically
reducing trainable parameters while preserving or exceeding the performance of full finetuning.

In addition, we introduce a \emph{Retrieval-Augmented Feature Enhancement} (RAFE) framework
that uses task-specific finetuned embeddings to retrieve nearest chemical neighbours from
ZINC-250k~\cite{gomez2018automatic}, a 249,455-molecule drug-like library, via GPU-accelerated
FAISS indices~\cite{johnson2021faiss}.
The physicochemical neighbourhood context derived from these retrievals (logP, QED, and SAS
distributions) provides an orthogonal, label-free signal that complements learned model
representations.

Our key contributions are:
\begin{itemize}
  \item A hybrid framework, \textbf{PEARL}, that supports three CLM families
        (ChemBERTa, MolFormer, Uni-Mol) with LoRA finetuning, achieving 75--96\% parameter
        reduction across all model sizes.
  \item A \textbf{3D geometry-aware} molecular property predictor via Uni-Mol LoRA, which
        yields the best BACE inhibitor classification performance (MCC~$= 0.623$) across all
        methods studied.
  \item A \textbf{Retrieval-Augmented Feature Enhancement} (RAFE) framework using ZINC-250k
        as an external, leakage-free knowledge base, establishing a new state of the art on
        ClinTox toxicity prediction (MCC~$= 0.913$) and matching end-to-end LoRA on BBBP.
  \item Comprehensive ablation across four molecular property prediction datasets, two loss
        functions, and three tree-based ensemble classifiers, demonstrating task-dependent
        optimal strategies.
\end{itemize}

All code is publicly available at \url{https://github.com/raghvendra5688/PEARL}.

%% ─────────────────────────────────────────────────────────────────────────────
\section{Materials and Methods}
%% ─────────────────────────────────────────────────────────────────────────────

\subsection{Data Collection and Curation}

\begin{itemize}
\item The \textbf{Blood-Brain Barrier Penetration (BBBP)} dataset was originally curated
  in~\cite{martins2012bayesian} for predicting barrier permeability and contains over 2,000
  compounds with binary labels, integrated into the MoleculeNet~\cite{wu2018moleculenet} benchmark.
\item \textbf{ClinTox} is part of the MoleculeNet benchmark~\cite{wu2018moleculenet}.
  The dataset contains 1,476 drugs that either failed clinical trials due to toxicity or succeeded,
  along with FDA approval status. FDA-approved drugs originate from the SWEETLEAD
  database~\cite{novick2013sweetlead}, while toxicity-failed candidates derive from the Aggregate
  Analysis of Clinical Trials (AACT).
\item The \textbf{Flavor Analysis and Recognition Transformer (FART)} dataset,
  introduced in~\cite{zimmermann2024fart}, is a multi-class taste prediction corpus of 15,025
  compounds annotated as \textit{bitter}, \textit{sour}, \textit{sweet}, \textit{umami},
  and \textit{undefined}.
\item The \textbf{BACE} dataset is part of MoleculeNet~\cite{wu2018moleculenet} and contains
  experimentally measured binding results for 1,513 inhibitors of human $\beta$-secretase~1
  (BACE-1), with binary activity labels. Dataset splits were curated using the DeepChem
  reference implementation~\cite{ramsundar2019deep}.
\end{itemize}

BBBP and ClinTox datasets are sourced from the IBM MolFormer repository and underwent
validity checks, duplicate removal, and class-ambiguity filtering.

\subsection{Data Splits and Class Distribution}

BBBP, ClinTox, and BACE employ scaffold-based splits~\cite{ramsundar2019deep} using an
80:10:10 train/validation/test ratio. BBBP exhibits moderate class imbalance (78.8\% permeable
in training). ClinTox faces severe imbalance: non-toxic compounds dominate at 92.3\% (CT\_TOX).
The Flavor dataset uses the curated splits provided in~\cite{zimmermann2024fart}, with
pronounced multi-class skew (sweet~41\%, bitter~32\%, sour~11\%, umami~4\%, undefined~12\%).
BACE is the most balanced, with 54.3\% inactive and 45.7\% active compounds.

\subsection{Pretrained Chemical Language Models}

\textbf{ChemBERTa-2}~\cite{ahmad2022chemberta2} builds upon
ChemBERTa~\cite{chithrananda2020chemberta} via a RoBERTa-style architecture trained on 77
million canonical SMILES from PubChem~\cite{kim2023pubchem}. Two pretraining strategies are
used: Masked Language Modeling (MLM) and Multi-Task Regression (MTR) on 200 molecular
property tasks. We use the two 77M-parameter variants, denoted \textbf{ChemBERTa-77M~MLM}
and \textbf{ChemBERTa-77M~MTR}, producing 384-dimensional embeddings.

\textbf{MolFormer}~\cite{ross2022molformer} is a scalable CLM that learns molecular
representations from SMILES without 3D geometry. It uses a transformer with linear attention
(FAVOR+)~\cite{likhosherstov2023favor}, rotary positional embeddings~\cite{heo2024rotary},
and efficient distributed training~\cite{isard2009distributed}, enabling pretraining on 1.1
billion molecules from PubChem~\cite{kim2023pubchem} and ZINC20~\cite{irwin2020zinc20}.
We use the \textbf{MolFormer-XL-10pct-both} variant ($\approx 46$M parameters, pretrained on
10\% of both datasets), producing 768-dimensional embeddings.

\textbf{Uni-Mol}~\cite{zhou2023unimol} is a universal 3D molecular pre-trained model
based on a SE(3)-equivariant Transformer with 15 attention layers.
Unlike SMILES-based CLMs, Uni-Mol accepts \emph{explicit 3D atomic coordinates} generated
via RDKit ETKDGv3 conformer embedding followed by MMFF94 force-field minimisation.
This geometry-aware design gives Uni-Mol inductive biases suited to structure-dependent
properties such as receptor binding affinity.
For PEARL, each molecule is represented as a
\textbf{2560-dimensional hybrid vector}: a 512-dimensional CLS token from the SE(3)-Transformer
concatenated with 2048-dimensional Morgan ECFP4 fingerprints (radius $= 2$), combining
3D structural awareness with explicit 2D topological information.

\subsection{Low-Rank Adaptation (LoRA)}
\label{sec:lora}

Low-Rank Adaptation (LoRA)~\cite{hu2021lora} introduces two trainable low-rank matrices $\mathbf{A} \in \mathbb{R}^{d \times r}$ and $\mathbf{B} \in \mathbb{R}^{r \times d}$ alongside
frozen pretrained weight matrices $\mathbf{W} \in \mathbb{R}^{d \times d}$ in the attention
layers, where $r \ll d$ is the rank hyperparameter. The weight update is:
\begin{equation}
  \mathbf{W}' = \mathbf{W} + \Delta\mathbf{W}, \quad \Delta\mathbf{W} = \mathbf{A} \cdot \mathbf{B}.
\end{equation}
Matrix $\mathbf{A}$ is initialised from $\mathcal{N}(0, \sigma^2)$ and $\mathbf{B}$ from
zeros, ensuring that adaptation begins from the pretrained representation.

\paragraph{LoRA for ChemBERTa and MolFormer.}
LoRA adapters are applied to the query ($\mathbf{W}^Q$), key ($\mathbf{W}^K$), and value
($\mathbf{W}^V$) projection matrices of each attention layer. All other parameters are frozen.
Hyperparameter search covers rank $r \in \{4, 8, 16, 32, 64\}$, LoRA scaling $\alpha \in \{8, 16, 32, 64, 128\}$, and dropout $\in \{0.0, 0.1, 0.2\}$.

\paragraph{LoRA for Uni-Mol.}
The version of \texttt{unimol\_tools} used in this work fuses $\mathbf{W}^Q$, $\mathbf{W}^K$,
and $\mathbf{W}^V$ into a single \texttt{in\_proj} weight matrix. LoRA adapters therefore
target \texttt{in\_proj} in every attention layer, with rank $r \in \{4, 8, 16, 32\}$ and
$\alpha \in \{8, 16, 32, 64\}$. To preserve gradient flow through the LoRA adapters, the
forward pass calls \texttt{\_repr.model.forward(\ldots, return\_repr=True)} directly rather
than the inference-mode \texttt{get\_repr()} wrapper.
Conformers for all dataset molecules are generated once before the hyperparameter sweep and
cached as \texttt{.pkl} files, avoiding redundant ETKDGv3 calls across sweep trials.
The \texttt{useExpTorsionAnglePrefs = False} flag prevents conformer generation failures on
metal-coordinated SMILES.

\paragraph{Training details.}
All models use AdamW with cosine-annealing learning rate schedules. The encoder (LoRA adapters)
receives the base learning rate; the classification head receives $5\times$ the base rate.
Gradient clipping ($\text{max\_norm} = 1.0$) is applied over all parameters. Early stopping
monitors validation MCC with patience of 5 epochs. Bayesian hyperparameter optimisation with
WandB sweeps (20--30 trials per dataset) optimises validation MCC.

\subsection{Loss Functions for Class Imbalance}

Two loss functions address the severe class imbalance prevalent in molecular datasets,
implemented via a custom extension of the HuggingFace \texttt{Trainer}:

\textbf{Weighted cross-entropy (WL):}
$\mathcal{L}_{\text{WL}} = -\sum_{i=0}^{C-1} w_i y_i \log\hat{y}_i$,
where $w_i = 1 - f_i$ (inverse class frequency). This gives the minority class higher loss
contribution and is most effective for severely imbalanced datasets (BBBP, Flavor, ClinTox).

\textbf{Focal loss (FL):}
$\mathcal{L}_{\text{FL}} = -\alpha (1 - \hat{p}_i)^\gamma \log\hat{p}_i$ with $\gamma = 2$,
which down-weights well-classified examples to focus learning on hard-to-classify boundary
cases. This excels on relatively balanced but challenging tasks such as BACE.

\subsection{PEARL Architecture}

PEARL defines four operational modes that form a progressive evaluation ladder:

\textbf{(a)~End-to-End (E2E) LoRA finetuning.}
The CLM backbone (ChemBERTa, MolFormer, or Uni-Mol) with LoRA adapters and a task-specific
classification head is optimised end-to-end for each downstream task.
For binary tasks, the head is a dense layer with sigmoid activation;
for multi-class tasks, a dense layer with softmax.

\textbf{(b)~Finetuned Embedding Prediction (FT~Embed).}
Task-specific embeddings are extracted from the finetuned model.
For MolFormer, the CLS token passes through two dense layers with GELU activation (768-dim
output). For ChemBERTa, the CLS token passes through a dense layer with tanh activation
(384-dim). For Uni-Mol, the CLS token is concatenated with ECFP4 fingerprints (2560-dim).
These embeddings train Optuna-tuned tree-based classifiers (XGBoost / LightGBM / CatBoost).

\textbf{(c)~FT~Embed $+$ PC (Physicochemical Feature Fusion).}
The FT~Embed vector is concatenated with 473~domain-specific physicochemical (PC) features:
148 RDKit molecular descriptors, 30 graph-based features (NetworkX v3.5), 128-bit Morgan
fingerprints, and 167-bit MACCS keys (RDKit v2025.09.1). This combined vector trains the
same tree-based classifiers without any retrieval from external databases.
A full description of engineered features is provided in Supplementary Table~S1.

\textbf{(d)~Retrieval-Augmented Feature Enhancement (RAFE).}
RAFE augments the FT~Embed~$+$~PC feature vector with ZINC-250k chemical neighbourhood
features ($\sim$73-dim; see Section~\ref{sec:rafe}), yielding the full feature vector
$[\text{CLM embeddings} \mid \text{PC features (473-dim)} \mid \text{ZINC-250k RAFE features}
(\sim\!73\text{-dim})]$, passed to the same Optuna-tuned tree-based classifiers.

\subsection{Retrieval-Augmented Feature Enhancement (RAFE)}
\label{sec:rafe}

\subsubsection{Design Rationale and External Knowledge Base}
\label{sec:zinc_rationale}

Naive retrieval-augmented approaches retrieve nearest neighbours from the task-specific training
set and use their labels as auxiliary features. This introduces three problems: a small retrieval
pool ($\sim$1,200--12,000 molecules), training-set leakage, and circular dependencies between
the retrieval space and the prediction target. PEARL addresses all three by using
\textbf{ZINC-250k}~\cite{gomez2018automatic} as the retrieval corpus---a curated subset of
249,455 drug-like molecules from the ZINC database~\cite{irwin2020zinc20} with no overlap with
any training, validation, or test split. Each ZINC-250k molecule carries three physicochemical
properties: lipophilicity (logP), Quantitative Estimate of Drug-likeness (QED, $\in [0,1]$),
and Synthetic Accessibility Score (SAS, $\in [1,10]$).

ZINC-250k is chosen as the retrieval corpus for four complementary reasons:

\begin{enumerate}
  \item \textbf{Scale and diversity.}
    249,455 drug-like molecules span a broad region of chemical space, providing a dense
    and representative neighbourhood for any query molecule from BACE, BBBP, ClinTox, or
    Flavor---datasets with at most 12,000 compounds.

  \item \textbf{Zero leakage.}
    ZINC-250k has no label overlap with any training, validation, or test split used in this
    study. Retrieved neighbours cannot carry prediction targets, eliminating the circular
    dependency that plagues training-set retrieval.

  \item \textbf{Label-free property signal.}
    logP, QED, and SAS are computed directly from SMILES without biological assay data.
    Their distributions in the local neighbourhood of a query molecule characterise its
    physicochemical environment in a way that is complementary to the molecule's own
    PC features and to the CLM's learned representation.

  \item \textbf{Cross-space consistency.}
    Embedding ZINC-250k independently with each finetuned CLM variant (ChemBERTa, MolFormer,
    Uni-Mol) enables cross-model neighbourhood comparison. Disagreement between ChemBERTa-space
    and MolFormer-space ZINC neighbourhoods is itself informative: molecules that look similar
    in one representation space but dissimilar in another are likely at the boundary of
    task-relevant chemical space.
\end{enumerate}

The core insight is that the logP/QED/SAS \emph{distributions of retrieved nearest ZINC
neighbours} characterise the local chemical space around a query molecule in a label-free,
leakage-free manner that complements the molecule's own PC features:
(i)~BBB-permeable molecules cluster in high-logP, high-QED ZINC regions~(BBBP);
(ii)~BACE inhibitors reside in drug-like, easily-synthesisable ZINC neighbourhoods~(BACE);
(iii)~clinically toxic compounds often exhibit unusual logP profiles and heterogeneous
physicochemical environments~(ClinTox);
(iv)~bitter compounds cluster in hydrophobic ZINC regions, while sweet/sour
compounds occupy hydrophilic regions~(Flavor).

\subsubsection{Vector Indexing with FAISS}

ZINC-250k molecules are embedded using task-specific finetuned model variants across all three
CLM families: six SMILES-based variants (ChemBERTa-MLM $\times$ 2 loss,
ChemBERTa-MTR $\times$ 2 loss, MolFormer $\times$ 2 loss) and two Uni-Mol variants
(FL and WL), yielding 32 embedding matrices in total (8 models $\times$ 4 datasets).
SMILES-based embeddings are 768-dimensional CLS tokens; Uni-Mol embeddings are
2560-dimensional hybrid vectors (512-dim SE(3)-Transformer CLS concatenated with
2048-dim ECFP4 fingerprints).
Embeddings are L2-normalised and indexed with FAISS
\texttt{GpuIndexFlatIP}~\cite{johnson2021faiss}---a GPU-accelerated exact cosine-similarity
search that achieves sub-2~ms query latency at 249k$\times$768 on a single GPU.
Indices are serialised as CPU \texttt{IndexFlatIP} files for portable, server-free retrieval.

\subsubsection{Chemical Neighbourhood Feature Extraction}

For every molecule in every split (train, validation, test), each FAISS index is queried
with $k \in \{10, 20\}$ nearest ZINC-250k neighbours. The following feature groups are
extracted:

\textbf{(a) Physicochemical neighbourhood profile ($\sim$27 features).}
Mean and standard deviation of logP, QED, and SAS across $k = 10$ and $k = 20$ neighbours,
plus similarity-weighted means (e.g.\ $\overline{\text{logP}}_w = \sum_i s_i \cdot
\text{logP}_i / \sum_i s_i$ where $s_i$ is cosine similarity to the $i$-th neighbour).

\textbf{(b) Neighbourhood density and isolation (4 features).}
Cosine similarity to the nearest ZINC neighbour ($\texttt{zinc\_nearest\_sim}$),
mean and standard deviation of similarities across $k = 10$ neighbours, and Shannon entropy
of the quantised similarity distribution. Molecules with low \texttt{zinc\_nearest\_sim}
are isolated in chemical space and their predictions are inherently less certain.

\textbf{(c) Embedding centroid features ($\sim$35 features).}
The mean of $k = 10$ neighbour embeddings forms a centroid; PCA fitted on training-split
centroids compresses this to 32 dimensions. Additional features capture cosine similarity of
the query to its centroid and the spread (mean L2 distance of neighbours from the centroid).

\textbf{(d) Tanimoto structural similarity (3 features).}
Tanimoto coefficients of Morgan fingerprints (radius~$= 2$, 128-bit) between the query and
its top-1 and top-5 ZINC neighbours, plus the fraction of top-10 neighbours sharing the same
Murcko scaffold.

\textbf{(e) Cross-space consistency (4 features).}
Each query produces two independent ZINC neighbourhood descriptions---one from ChemBERTa
embedding space and one from MolFormer space. The four features capture disagreement in
weighted-logP, weighted-QED, Jaccard overlap of the two $k = 10$ neighbour sets, and the
ratio of neighbourhood densities between the two spaces.

The full augmented feature vector per molecule is:
$[\text{CLM embeddings} \mid \text{PC features (473-dim)} \mid \text{ZINC RAFE features}
(\sim\!73\text{-dim})]$.

\subsection{Tree-based Ensemble Classifiers}

The final prediction stage employs XGBoost~\cite{chen2016xgboost}, LightGBM~\cite{ke2017lightgbm}, and CatBoost~\cite{prokhorenkova2019catboost}.
These algorithms handle high-dimensional mixed-type feature vectors, provide built-in feature
importance, and are robust to outliers. Bayesian hyperparameter optimisation with
Optuna~\cite{akiba2019optuna} (10 trials for base configurations, 20 for RAFE) maximises
validation MCC via stratified 5-fold cross-validation. Class imbalance is handled via
\texttt{scale\_pos\_weight} (XGBoost), \texttt{class\_weight="balanced"} (LightGBM), and
\texttt{auto\_class\_weights="Balanced"} (CatBoost).

\subsection{Evaluation Metrics}

Performance is evaluated using accuracy (Acc), area under the ROC curve (AUC), Matthews
correlation coefficient (MCC), precision (Prec), recall (Rec), and F1-macro (F1$_\text{mac}$)
/ F1-micro (F1$_\text{mic}$). MCC is the primary optimisation target due to its robustness
to class imbalance. Area under the precision-recall curve (AUPR) is additionally reported.

%% ─────────────────────────────────────────────────────────────────────────────
\section{PEARL Architecture Overview}
%% ─────────────────────────────────────────────────────────────────────────────

Figure~\ref{fig:architecture} illustrates the complete PEARL workflow. The SMILES
string of each molecule is processed by one of three CLM backbones.
LoRA adapters augment only the attention projection layers, keeping all other pretrained
weights frozen.
For Uni-Mol, 3D conformers are pre-generated via ETKDGv3 + MMFF94 and cached.
The finetuned model provides both end-to-end predictions and task-specific embeddings.
These embeddings are optionally concatenated with engineered PC features and with ZINC-250k
RAFE features retrieved via FAISS indices, and the resulting feature vector is fed to Optuna-tuned
tree-based classifiers.

\begin{figure}[t]
  \centering
  \includegraphics[width=\linewidth]{figures/effichm2_architecture.pdf}
  \caption{PEARL architecture. A query SMILES is processed by one of three CLM
  families (ChemBERTa, MolFormer, Uni-Mol) with LoRA adapters. Task-specific embeddings
  are fused with molecular PC features and ZINC-250k retrieval features, then passed to
  tree-based classifiers. Uni-Mol additionally encodes 3D ETKDGv3 conformers.}
  \label{fig:architecture}
\end{figure}

%% ─────────────────────────────────────────────────────────────────────────────
\section{Experimental Results}
%% ─────────────────────────────────────────────────────────────────────────────

Results are organised around the four PEARL operational modes described in Section~2.4:
(i)~\textbf{E2E LoRA}---end-to-end LoRA finetuning of the CLM backbone;
(ii)~\textbf{FT Embed}---Optuna-tuned tree classifiers on frozen finetuned embeddings;
(iii)~\textbf{FT Embed $+$ PC}---FT Embed augmented with 473-dim physicochemical features,
without ZINC-250k retrieval;
(iv)~\textbf{RAFE}---FT Embed $+$ PC further augmented with ZINC-250k neighbourhood features.
For each mode we report the best-performing CLM variant and, where applicable, tree model.
Full ablation tables appear in the Supplementary Material.

\subsection{Sequence of Comparisons}
\label{sec:comparison_sequence}

The four PEARL modes form a deliberate evaluation ladder that isolates the contribution
of each component:

\textbf{(i)~E2E LoRA establishes the performance ceiling} achievable by fully differentiable
gradient flow through the CLM backbone.  Because LoRA adapters and the classification head
are optimised jointly, the model can learn task-specific representations end-to-end.
E2E LoRA results are therefore treated as the upper bound for each CLM family and dataset.

\textbf{(ii)~FT Embed quantifies how much embedding quality alone contributes.}
Here, the finetuned model is frozen after training; its CLS token embedding is extracted and
passed to an Optuna-tuned tree classifier (XGBoost / LightGBM / CatBoost).
Any gap between FT Embed and E2E LoRA reflects information that can only be recovered through
end-to-end gradient flow---i.e., representations that emerge from joint task-head optimisation
but are not fully preserved in the frozen CLS token.

\textbf{(iii)~FT Embed $+$ PC isolates the contribution of physicochemical feature engineering.}
The frozen embedding is augmented with 473 engineered PC features and passed to the same
tree classifier. A gain over FT Embed reflects the additional discriminative power of
handcrafted molecular descriptors and fingerprints, independent of any external retrieval.

\textbf{(iv)~RAFE measures the additional value of ZINC-250k chemical neighbourhood context.}
RAFE augments the FT Embed $+$ PC feature vector with ZINC-250k retrieval features
($\sim$73-dim), all passed to the same tree classifier family.
A gain from FT Embed $+$ PC to RAFE indicates that neighbourhood signals from ZINC-250k
provide information orthogonal to both the learned CLM representation and the molecule's own
physicochemical profile.
A lack of gain (or degradation) indicates that the task is primarily driven by geometry-critical
or training-distribution-specific signals that retrieval cannot replicate.

This ladder allows us to attribute performance differences to four distinct sources:
(a)~the quality of the underlying pre-trained CLM, (b)~the fidelity of the frozen embedding,
(c)~the discriminative power of engineered PC features, and
(d)~the neighbourhood context contributed by ZINC-250k.
The per-dataset results in the following subsections are discussed in this light.

\subsection{BBBP Dataset Results}

Table~\ref{tab:bbbp} presents the best configuration per modality and CLM family for BBBP.
E2E LoRA finetuning of MolFormer-XL with weighted loss achieves optimal overall performance
(MCC~$= 0.797$, AUC~$= 0.948$, F1$_\text{mic} = 0.906$), matching the best previously reported result for this dataset.

Uni-Mol with weighted loss achieves MCC~$= 0.774$ and AUC~$= 0.943$, demonstrating that
3D geometry-aware representations are highly competitive on BBBP, despite this task being
conventionally attributed to 2D physicochemical properties such as logP and molecular weight.
Uni-Mol FT Embed (WL + CatBoost) achieves MCC~$= 0.720$, AUC~$= 0.918$---competitive
with MolFormer FT Embed while using 3D-geometry-derived embeddings.
Applying RAFE to Uni-Mol embeddings (FL + XGBoost) improves this to MCC~$= 0.773$,
AUC~$= 0.928$---marginally below Uni-Mol E2E LoRA (MCC~$= 0.774$) while using only
tree-based classifiers, and narrowing the gap with MolFormer RAFE (MCC~$= 0.798$).

The RAFE framework (MolFormer-XL FL + XGBoost) achieves MCC~$= 0.798$---\emph{matching} the
E2E LoRA result while using a tree-based classifier rather than end-to-end neural finetuning.
Importantly, RAFE provides an interpretable neighbourhood context signal: the high mean logP
of retrieved ZINC-250k neighbours (associated with BBB-permeable molecules) directly explains
the prediction.
Adding PC features to finetuned embeddings without RAFE (PC + FT Embed,
MolFormer WL + CatBoost: MCC~$= 0.773$) is outperformed by both E2E LoRA and RAFE,
suggesting that ZINC-250k neighbourhood context contributes information orthogonal to the
molecule's own PC features.

\begin{table*}[t]
\centering
\caption{Best performing configurations for each modality and CLM on the BBBP dataset.
$^*$ denotes the top-performing method by average metric. RAFE = FT~Embed~+~PC~+~ZINC-250k
retrieval features.}
\label{tab:bbbp}
\resizebox{\textwidth}{!}{%
\begin{tabular}{lllllrrrrrrr}
\toprule
Modality & CLM & Loss & ML & AUC & AUPR & Acc & F1$_\text{mac}$ & F1$_\text{mic}$ & MCC & Prec & Rec \\
\midrule
E2E LoRA & ChemBERTa-77M MLM & FL & --- & 0.923 & --- & 0.880 & 0.851 & 0.865 & 0.739 & 0.889 & 0.851 \\
E2E LoRA & ChemBERTa-77M MTR & FL & --- & 0.921 & --- & 0.870 & 0.843 & 0.854 & 0.715 & 0.873 & 0.843 \\
E2E LoRA$^*$ & MolFormer-XL & WL & --- & 0.948 & --- & 0.906 & 0.898 & 0.906 & 0.797 & 0.901 & 0.896 \\
E2E LoRA & Uni-Mol & FL & --- & 0.944 & 0.949 & 0.880 & 0.873 & 0.880 & 0.747 & 0.868 & 0.878 \\
E2E LoRA & Uni-Mol & WL & --- & 0.943 & 0.955 & 0.896 & 0.887 & 0.896 & 0.774 & 0.890 & 0.885 \\
\midrule
FT Embed & ChemBERTa-77M MLM & FL & LightGBM & 0.905 & --- & 0.896 & 0.882 & 0.896 & 0.776 & 0.911 & 0.866 \\
FT Embed & MolFormer-XL & FL & XGBoost & 0.935 & --- & 0.901 & 0.888 & 0.901 & 0.787 & 0.914 & 0.873 \\
FT Embed & Uni-Mol & FL & XGBoost & 0.917 & 0.939 & 0.865 & 0.842 & 0.865 & 0.709 & 0.889 & 0.823 \\
FT Embed & Uni-Mol & WL & CatBoost & 0.918 & 0.917 & 0.865 & 0.838 & 0.865 & 0.720 & 0.912 & 0.814 \\
\midrule
PC + FT Embed & ChemBERTa-77M MLM & FL & LightGBM & 0.888 & --- & 0.891 & 0.875 & 0.891 & 0.765 & 0.907 & 0.859 \\
PC + FT Embed & MolFormer-XL & FL & LightGBM & 0.934 & --- & 0.901 & 0.887 & 0.901 & 0.789 & 0.920 & 0.870 \\
PC + FT Embed & MolFormer-XL & WL & CatBoost & 0.951 & --- & 0.896 & 0.884 & 0.896 & 0.773 & 0.902 & 0.872 \\
\midrule
RAFE & ChemBERTa-77M MLM & FL & LightGBM & 0.926 & 0.942 & 0.896 & 0.882 & 0.896 & 0.776 & 0.911 & 0.866 \\
RAFE & MolFormer-XL & FL & XGBoost & \textbf{0.940} & 0.940 & \textbf{0.906} & \textbf{0.895} & \textbf{0.906} & \textbf{0.798} & 0.918 & 0.881 \\
RAFE & MolFormer-XL & WL & CatBoost & 0.946 & \textbf{0.960} & 0.885 & 0.872 & 0.885 & 0.750 & 0.889 & 0.861 \\
RAFE & Uni-Mol & FL & XGBoost & 0.928 & 0.943 & 0.891 & 0.872 & 0.891 & 0.773 & 0.927 & 0.850 \\
\bottomrule
\end{tabular}
}
\end{table*}

\subsection{Flavor Dataset Results}

Table~\ref{tab:flavor} presents results on the multi-class Flavor dataset.
E2E LoRA finetuning of MolFormer-XL with weighted loss achieves MCC~$= 0.801$,
F1$_\text{mac} = 0.809$, and AUC~$= 0.975$, with balanced precision and recall.

The RAFE framework (MolFormer-XL WL + CatBoost) achieves MCC~$= 0.804$, marginally
\emph{exceeding} E2E LoRA, while using tree-based classifiers on the augmented feature
vector. This suggests that ZINC-250k neighbourhood context contributes orthogonal
information for distinguishing multi-class taste profiles: logP differences between the
hydrophobic ZINC neighbourhood of bitter compounds versus the hydrophilic neighbourhood
of sweet/sour compounds complement the CLM's learned representations.

Uni-Mol achieves MCC~$= 0.776$ with weighted loss, competitive with but below MolFormer,
indicating that 3D geometric information does not provide a strong additional signal for
flavour prediction compared to SMILES-based representations.
Uni-Mol FT Embed FL (XGBoost) achieves MCC~$= 0.765$, Acc~$= 0.866$, F1$_\text{mac} = 0.731$,
matching Uni-Mol E2E LoRA WL (MCC~$= 0.776$) closely and slightly outperforming E2E LoRA FL
(MCC~$= 0.765$), demonstrating that finetuned 3D embeddings alone capture competitive taste
representations without end-to-end gradient flow.
Uni-Mol RAFE FL (CatBoost) achieves MCC~$= 0.764$, Acc~$= 0.866$, F1$_\text{mac} = 0.714$,
while WL (CatBoost) achieves MCC~$= 0.761$, F1$_\text{mac} = 0.737$---WL's minority-class
upweighting yields superior macro-averaged balance, particularly for the rarest umami class
(AUPR: WL $= 0.507$ vs FL $= 0.379$), while FL edges ahead on overall MCC.

\begin{table*}[t]
\centering
\caption{Best performing configurations on the Flavor (FART) dataset. RAFE results include
per-class AUPR for the five taste categories: bitter (C0), sweet (C1), sour (C2),
umami (C3), undefined (C4).}
\label{tab:flavor}
\resizebox{\textwidth}{!}{%
\begin{tabular}{llllrrrrr}
\toprule
Modality & CLM & Loss & ML & AUC & Acc & F1$_\text{mac}$ & F1$_\text{mic}$ & MCC \\
\midrule
E2E LoRA$^*$ & MolFormer-XL & WL & --- & \textbf{0.975} & 0.892 & 0.809 & 0.892 & 0.801 \\
E2E LoRA & ChemBERTa-77M MTR & WL & --- & 0.964 & 0.890 & 0.780 & 0.890 & 0.802 \\
E2E LoRA & Uni-Mol & FL & --- & 0.966 & 0.868 & 0.770 & 0.868 & 0.765 \\
E2E LoRA & Uni-Mol & WL & --- & 0.967 & 0.874 & 0.772 & 0.874 & 0.776 \\
\midrule
FT Embed & MolFormer-XL & WL & LightGBM & 0.970 & 0.882 & 0.783 & 0.882 & 0.787 \\
FT Embed & Uni-Mol & FL & XGBoost & 0.950 & 0.866 & 0.731 & 0.866 & 0.765 \\
PC + FT Embed & MolFormer-XL & WL & LightGBM & 0.967 & 0.883 & 0.798 & 0.883 & 0.789 \\
\midrule
RAFE & ChemBERTa-77M MTR & FL & CatBoost & 0.824 & 0.886 & 0.741 & 0.886 & 0.798 \\
RAFE & MolFormer-XL & FL & LightGBM & 0.811 & 0.889 & 0.788 & 0.889 & 0.793 \\
RAFE & MolFormer-XL & WL & LightGBM & 0.831 & 0.892 & \textbf{0.794} & 0.892 & 0.798 \\
RAFE & MolFormer-XL & WL & CatBoost & 0.837 & \textbf{0.893} & 0.774 & \textbf{0.893} & \textbf{0.804} \\
RAFE & Uni-Mol & FL & CatBoost & --- & 0.866 & 0.714 & 0.866 & 0.764 \\
\bottomrule
\end{tabular}
}
\end{table*}

\subsection{BACE Dataset Results}

Table~\ref{tab:bace} compares key configurations on the BACE inhibitor prediction task.
BACE is the most challenging dataset (MCC $< 0.65$ across most methods), consistent with the
difficulty of generalising to novel chemical scaffolds in a binding-affinity prediction task.

\textbf{Uni-Mol with weighted loss achieves the best overall performance} (MCC~$= 0.623$,
AUC~$= 0.891$, Acc~$= 0.822$), representing an improvement of 8\% in MCC over the previous
best E2E LoRA MolFormer FL (MCC~$= 0.577$, AUC~$= 0.872$).
This improvement is consistent with the hypothesis that 3D geometry-aware representations
provide superior inductive biases for binding affinity tasks: BACE-1 inhibitor activity is
driven by ligand--receptor complementarity in three-dimensional space.

The RAFE framework with MolFormer-XL WL + LightGBM achieves MCC~$= 0.561$, intermediate
between the finetuned-embedding-only result (MCC~$= 0.557$) and E2E LoRA (MCC~$= 0.577$).
Adding PC features to finetuned embeddings provides minimal improvement over RAFE alone,
consistent with the finding that LoRA-adapted representations already capture
task-relevant molecular properties.
Uni-Mol FT Embed (FL + LightGBM) achieves MCC~$= 0.493$, AUC~$= 0.860$---substantially
below Uni-Mol E2E LoRA (MCC~$= 0.623$)---confirming that end-to-end 3D gradient flow is
essential for binding affinity prediction and cannot be fully recovered at the embedding stage.
Uni-Mol RAFE (FL + XGBoost, MCC~$= 0.478$, AUC~$= 0.840$) modestly improves over
Uni-Mol FT Embed but remains well below Uni-Mol E2E LoRA, reaffirming that
geometry-critical tasks benefit most from fully differentiable 3D encoding.

Figure~\ref{fig:bace} shows ROC and precision-recall curves for the best configuration per
model family. MolFormer-XL base embeddings with CatBoost achieve AUC~$= 0.885$ (the highest
among SMILES-based approaches), while Uni-Mol WL achieves the highest AUPR~$= 0.924$ and
overall MCC.

\begin{table}[ht]
\centering
\caption{Best configurations on the BACE dataset. $^*$ overall best by MCC.}
\label{tab:bace}
\resizebox{\columnwidth}{!}{%
\begin{tabular}{llllrrrr}
\toprule
Modality & CLM & Loss & ML & AUC & AUPR & Acc & MCC \\
\midrule
E2E LoRA & MolFormer-XL & FL & --- & 0.872 & 0.902 & 0.789 & 0.577 \\
E2E LoRA & MolFormer-XL & WL & --- & 0.871 & 0.897 & 0.724 & 0.476 \\
E2E LoRA & Uni-Mol & FL & --- & 0.876 & 0.924 & 0.743 & 0.520 \\
E2E LoRA$^*$ & Uni-Mol & WL & --- & \textbf{0.891} & \textbf{0.924} & \textbf{0.822} & \textbf{0.623} \\
\midrule
FT Embed & MolFormer-XL & FL & CatBoost & 0.869 & 0.895 & 0.763 & 0.557 \\
FT Embed & Uni-Mol & FL & LightGBM & 0.860 & 0.893 & 0.730 & 0.493 \\
PC + FT Embed & MolFormer-XL & FL & XGBoost & 0.870 & 0.901 & 0.737 & 0.517 \\
\midrule
RAFE & MolFormer-XL & FL & CatBoost & 0.863 & 0.892 & 0.763 & 0.557 \\
RAFE & MolFormer-XL & WL & LightGBM & 0.858 & 0.892 & 0.770 & 0.561 \\
RAFE & Uni-Mol & FL & XGBoost & 0.840 & 0.895 & 0.704 & 0.478 \\
\bottomrule
\end{tabular}
}
\end{table}

\subsection{ClinTox Dataset Results}

Table~\ref{tab:clintox} presents results for the CT\_TOX (clinical toxicity) and FDA\_APPROVED
tasks.

The \textbf{RAFE framework with MolFormer-XL FL + CatBoost establishes a new state of the art}
on the CT\_TOX task: MCC~$= 0.913$, AUC~$= 0.986$, Acc~$= 0.986$, F1$_\text{mac} = 0.955$---
surpassing E2E LoRA (MolFormer FL, MCC~$= 0.786$) by 12.7\% and FT~Embed (MolFormer FL,
MCC~$= 0.827$) by 8.6\%, the largest RAFE gain across all four datasets.
This result is particularly notable because ClinTox uniquely benefits from ZINC-250k
neighbourhood context: the retrieval pipeline identifies atypical logP/SAS signatures
associated with clinically toxic compounds that are not captured by finetuned embeddings alone.

Uni-Mol E2E LoRA performs poorly on ClinTox (MCC~$\leq 0.282$), due to extreme class
imbalance ($>$90\% non-toxic) combined with the small dataset size, making it difficult to
learn a reliable 3D conformer-based toxicity signal end-to-end.
However, Uni-Mol FT Embed (FL + CatBoost) substantially recovers performance
(MCC~$= 0.460$, AUC~$= 0.923$, AUPR~$= 0.996$), and Uni-Mol RAFE with weighted loss
(WL + XGBoost) achieves MCC~$= 0.633$, AUC~$= 0.983$, AUPR~$= 0.999$,
Acc~$= 0.951$---demonstrating that ZINC-250k neighbourhood context compensates for E2E
training limitations under severe class imbalance, and that weighted-loss embeddings better
capture the physicochemical signatures of clinically toxic compounds for retrieval.

For FDA\_APPROVED, ChemBERTa-MLM WL + LightGBM via RAFE achieves
AUC~$= 0.992$, AUPR~$= 1.000$, and MCC~$= 0.827$.

\begin{table}[ht]
\centering
\caption{Best configurations on the ClinTox dataset. Results shown for CT\_TOX (toxicity)
and FDA\_APPROVED labels. $^*$ overall best by MCC.}
\label{tab:clintox}
\resizebox{\columnwidth}{!}{%
\begin{tabular}{llllllrrr}
\toprule
Task & Modality & CLM & Loss & ML & AUC & AUPR & Acc & MCC \\
\midrule
\multirow{8}{*}{CT\_TOX}
 & E2E LoRA & MolFormer-XL & FL & --- & 0.993 & --- & 0.972 & 0.786 \\
 & E2E LoRA & Uni-Mol & FL & --- & 0.904 & 0.995 & 0.867 & 0.282 \\
 & E2E LoRA & Uni-Mol & WL & --- & 0.802 & 0.988 & 0.937 & 0.276 \\
 & FT Embed & Uni-Mol & FL & CatBoost & 0.923 & 0.996 & 0.853 & 0.460 \\
 & RAFE & ChemBERTa-MLM & WL & LightGBM & 0.989 & 0.939 & 0.979 & 0.869 \\
 & RAFE & MolFormer-XL & FL & LightGBM & 0.997 & 0.975 & 0.979 & 0.867 \\
 & RAFE$^*$ & MolFormer-XL & FL & CatBoost & 0.986 & 0.931 & \textbf{0.986} & \textbf{0.913} \\
 & RAFE & Uni-Mol & WL & XGBoost & 0.983 & 0.999 & 0.951 & 0.633 \\
\midrule
\multirow{4}{*}{\small FDA\_APP.}
 & E2E LoRA & MolFormer-XL & FL & --- & 0.993 & --- & 0.972 & 0.786 \\
 & FT Embed & MolFormer-XL & FL & XGBoost & 0.987 & --- & 0.979 & 0.827 \\
 & RAFE & ChemBERTa-MLM & WL & LightGBM & 0.992 & 1.000 & 0.979 & 0.827 \\
 & RAFE & MolFormer-XL & FL & XGBoost & 0.988 & 0.999 & 0.979 & 0.827 \\
\bottomrule
\end{tabular}
}
\end{table}

\subsection{Retrieval-Augmented Feature Enhancement Analysis}

Table~\ref{tab:rafe_summary} provides a consolidated comparison of all four PEARL modes
across all four datasets: E2E LoRA, FT~Embed (frozen embeddings only), FT~Embed+PC
(embeddings with 473-dim physicochemical features, no ZINC-250k retrieval), and RAFE
(full retrieval-augmented pipeline), including Uni-Mol variants.

For BBBP, MolFormer RAFE (MCC~$= 0.798$) essentially ties E2E LoRA (MCC~$= 0.797$),
demonstrating that an interpretable tree-based approach with chemical neighbourhood context
can match neural end-to-end finetuning. Uni-Mol RAFE FL (MCC~$= 0.773$) is nearly equivalent
to Uni-Mol E2E LoRA (MCC~$= 0.774$) using only a tree-based classifier, and narrows the gap
with MolFormer RAFE substantially (from $\Delta = 0.024$ at E2E to $\Delta = 0.025$ at RAFE).
For Flavor, MolFormer RAFE marginally improves over E2E LoRA (MCC~$= 0.804$ vs.\ $0.801$),
while Uni-Mol FT Embed FL (XGBoost, MCC~$= 0.765$) closely matches Uni-Mol E2E LoRA FL
(MCC~$= 0.765$), confirming that finetuned 3D features transfer well to tree classifiers
for multi-class taste prediction.
Uni-Mol RAFE FL (MCC~$= 0.764$) falls marginally below FT Embed FL, confirming
that ZINC-250k retrieval provides no additional benefit for 3D geometry in taste prediction.
For BACE, Uni-Mol E2E (MCC~$= 0.623$) substantially outperforms all RAFE configurations
(MolFormer RAFE: $0.561$; Uni-Mol RAFE FL: $0.478$), confirming that geometry-critical tasks
require end-to-end 3D gradient flow.
For ClinTox CT\_TOX, MolFormer RAFE delivers the largest improvement over E2E LoRA (MCC~$= 0.913$ vs.\ $0.786$),
and Uni-Mol RAFE with weighted loss (WL + XGBoost, MCC~$= 0.633$) dramatically outperforms
Uni-Mol E2E LoRA (MCC~$= 0.282$), showing that weighted-loss embeddings better encode the
physicochemical signature of toxic compounds for ZINC-250k retrieval.

\begin{table}[ht]
\centering
\caption{Summary comparison of the four PEARL modes across all four datasets (best MCC per
approach). E2E LoRA results use the best CLM variant per dataset. FT~Embed and FT~Embed+PC
use MolFormer-XL unless otherwise noted. ``RAFE best (MolFormer)'' uses MolFormer-XL
embeddings. ``Uni-Mol RAFE'' uses 2560-dim Uni-Mol hybrid embeddings.}
\label{tab:rafe_summary}
\resizebox{\columnwidth}{!}{%
\begin{tabular}{lrrrrr}
\toprule
Dataset & Approach & AUC & Acc & F1$_\text{mac}$ & MCC \\
\midrule
\multirow{6}{*}{BBBP}
 & E2E LoRA (MolFormer WL)    & 0.948 & 0.906 & 0.898 & 0.797 \\
 & Uni-Mol E2E (WL)           & 0.943 & 0.896 & 0.887 & 0.774 \\
 & FT Embed (MolFormer FL)    & 0.935 & 0.901 & 0.888 & 0.787 \\
 & FT Embed+PC (MolFormer FL) & 0.934 & 0.901 & 0.887 & 0.789 \\
 & RAFE best (MolFormer)      & 0.940 & 0.906 & 0.895 & \textbf{0.798} \\
 & Uni-Mol RAFE (FL)          & 0.928 & 0.891 & 0.872 & 0.773 \\
\midrule
\multirow{6}{*}{Flavor}
 & E2E LoRA (MolFormer WL)    & 0.975 & 0.892 & 0.809 & 0.801 \\
 & Uni-Mol E2E (WL)           & 0.967 & 0.874 & 0.772 & 0.776 \\
 & FT Embed (MolFormer WL)    & 0.970 & 0.882 & 0.783 & 0.787 \\
 & FT Embed+PC (MolFormer WL) & 0.967 & 0.883 & 0.798 & 0.789 \\
 & RAFE best (MolFormer)      & 0.837 & 0.893 & 0.774 & \textbf{0.804} \\
 & Uni-Mol RAFE (FL)          & ---   & 0.866 & 0.714 & 0.764 \\
\midrule
\multirow{6}{*}{BACE}
 & E2E LoRA (MolFormer FL)    & 0.872 & 0.789 & 0.785 & 0.577 \\
 & Uni-Mol E2E (WL)           & 0.891 & 0.822 & 0.808 & \textbf{0.623} \\
 & FT Embed (MolFormer FL)    & 0.869 & 0.763 & ---   & 0.557 \\
 & FT Embed+PC (MolFormer FL) & 0.870 & 0.737 & ---   & 0.517 \\
 & RAFE best (MolFormer)      & 0.858 & 0.770 & 0.768 & 0.561 \\
 & Uni-Mol RAFE (FL)          & 0.840 & 0.704 & 0.704 & 0.478 \\
\midrule
\multirow{4}{*}{ClinTox}
 & Uni-Mol E2E (FL)        & 0.904 & 0.867 & 0.612 & 0.282 \\
 & FT Embed best (MolFormer) & 0.987 & 0.979 & 0.906 & 0.827 \\
 & RAFE best (MolFormer)   & 0.986 & 0.986 & 0.955 & \textbf{0.913} \\
 & Uni-Mol RAFE (WL)       & 0.983 & 0.951 & 0.803 & 0.633 \\
\bottomrule
\end{tabular}
}
\end{table}

\subsection{Parameter Efficiency}

Table~\ref{tab:params} reports LoRA hyperparameters and trainable parameter counts for all
model families across the four tasks. ChemBERTa variants (3.4M base parameters) achieve
reductions of 76--94\% depending on rank; MolFormer-XL ($\approx 46$M parameters) achieves
84--97\% reduction.
For Uni-Mol (rank $r \in \{4,8,16,32\}$), LoRA targets only the \texttt{in\_proj}
weight in 15 attention layers, yielding a particularly efficient adaptation with $<5\%$
trainable parameters for the lowest-rank configurations.

Full parameter counts are tabulated in Supplementary Table~S6.

\begin{table}[ht]
\centering
\caption{Selected LoRA configurations illustrating the range of parameter efficiency.
Trainable~\% = 100 $\times$ trainable / total parameters.}
\label{tab:params}
\resizebox{\columnwidth}{!}{%
\begin{tabular}{llllrrrr}
\toprule
Dataset & Model & $r$ & $\alpha$ & Total & Trainable & \% & \% Red. \\
\midrule
BACE   & ChemBERTa-MLM & 4  & 32  & 3.43M & 205K  & 5.66  & 94.34 \\
BACE   & ChemBERTa-MTR & 32 & 128 & 3.43M & 606K  & 15.03 & 84.97 \\
BACE   & MolFormer-XL  & 64 & 128 & 45.6M & 8.26M & 15.35 & 84.65 \\
BBBP   & ChemBERTa-MLM & 64 & 32  & 3.43M & 1.06M & 23.69 & 76.31 \\
BBBP   & MolFormer-XL  & 16 & 16  & 45.6M & 2.95M & 6.09  & 93.91 \\
ClinTox& MolFormer-XL  & 4  & 32  & 45.6M & 1.63M & 3.44  & 96.56 \\
Flavor & MolFormer-XL  & 4  & 64  & 45.6M & 1.63M & 3.44  & 96.56 \\
\bottomrule
\end{tabular}
}
\end{table}

%% ─────────────────────────────────────────────────────────────────────────────
\section{Discussion \& Conclusion}
%% ─────────────────────────────────────────────────────────────────────────────

PEARL extends parameter-efficient molecular property prediction in two directions that
address limitations of SMILES-only, single-modality CLM finetuning.

\textbf{3D geometry matters for binding-affinity tasks.}
Uni-Mol's SE(3)-equivariant architecture provides the most accurate BACE inhibitor
classification (MCC~$= 0.623$ vs.\ $0.577$ for the best SMILES-based approach).
BACE-1 inhibitor activity is governed by three-dimensional ligand--receptor complementarity,
and the inductive biases of a geometry-aware transformer are therefore well-matched to this
task. The hybrid 2560-dimensional representation (512-dim 3D CLS token plus 2048-dim ECFP4)
balances 3D structural awareness with explicit 2D topological information. Conversely,
Uni-Mol does not improve on SMILES-based models for ClinTox, where the primary challenge
is severe class imbalance rather than geometric discrimination.

\textbf{Retrieval-augmented chemical neighbourhood context provides orthogonal information.}
The RAFE framework demonstrates that ZINC-250k neighbourhood features---logP/QED/SAS
distributions, neighbourhood density, cross-space consistency---provide information
complementary to learned model representations. On ClinTox CT\_TOX, where finetuning
paradoxically \emph{degrades} performance in the non-RAFE setting (likely due to overfitting
under extreme class imbalance), RAFE restores and surpasses the original best, achieving
MCC~$= 0.913$. This suggests that the ZINC-250k context provides a regulatory signal that
helps the model identify the unusual physicochemical profiles of clinical toxic compounds
without overfitting to the small labelled dataset. On BBBP, RAFE achieves MCC~$= 0.798$,
essentially matching E2E LoRA (MCC~$= 0.797$) with a fully interpretable tree-based pipeline.

\textbf{Task-dependent optimal strategies persist.}
Several consistent patterns emerge across the expanded study:
(i)~MolFormer consistently outperforms ChemBERTa across all experimental conditions and tasks;
(ii)~weighted loss excels on severely imbalanced datasets (BBBP, Flavor, ClinTox), while focal
loss is optimal for challenging but relatively balanced tasks (BACE);
(iii)~adding PC features to finetuned embeddings provides minimal improvement over finetuned
embeddings alone, because LoRA adaptation effectively internalises task-relevant molecular
properties;
(iv)~for ClinTox, base embeddings combined with PC features remain a strong baseline, and
RAFE is needed to outperform them with finetuned representations.

\textbf{Parameter efficiency.}
LoRA achieves 75--97\% parameter reduction across all three model families without sacrificing
performance. The smallest effective configurations (MolFormer-XL at $r = 4$, $\alpha = 64$
for Flavor and ClinTox, 3.44\% trainable) are particularly noteworthy: they demonstrate that
the intrinsic task-relevant subspace of a 46M-parameter molecular transformer can be captured
with fewer than 1.7M additional parameters.

\textbf{Uni-Mol RAFE: geometry-aware retrieval amplifies performance selectively.}
Applying RAFE to Uni-Mol 2560-dim embeddings reveals a consistent, task-dependent pattern,
with loss function choice critically influencing retrieval quality.
For BBBP, Uni-Mol RAFE FL (MCC~$= 0.773$) nearly matches Uni-Mol E2E LoRA (MCC~$= 0.774$)
using only a tree-based classifier, and narrows the gap with MolFormer RAFE (MCC~$= 0.798$)---
focal-loss embeddings retrieve ZINC-250k neighbours with the highest-logP physicochemical
profiles characteristic of BBB-permeable compounds.
For BACE---the geometry-critical task---Uni-Mol E2E LoRA (MCC~$= 0.623$) remains
substantially superior to both Uni-Mol FT Embed (MCC~$= 0.493$) and Uni-Mol RAFE FL
(MCC~$= 0.478$), confirming that end-to-end 3D gradient flow is irreplaceable for
binding affinity prediction and cannot be recovered through retrieval.
For ClinTox, Uni-Mol RAFE \emph{with weighted loss} (WL + XGBoost, MCC~$= 0.633$, AUC~$= 0.983$)
dramatically outperforms both Uni-Mol E2E LoRA (MCC~$= 0.282$) and the FL variant
(best MCC~$= 0.547$)---weighted-loss training upweights the minority toxic class during
finetuning, producing embeddings that retrieve ZINC-250k neighbours with atypical logP/SAS
profiles characteristic of clinically toxic compounds, providing a reliable surrogate signal
where end-to-end 3D learning fails under extreme class imbalance.

\textbf{Future directions.}
Several avenues remain open.
First, extending PEARL to regression tasks (binding affinity, solubility) with
calibrated uncertainty estimates would broaden its applicability.
Second, emerging multi-modal CLMs such as nach0~\cite{livne2024nach0} and
Omni-Mol~\cite{hu2024omnimol}, which integrate SMILES, graphs, and textual descriptions,
represent natural extensions of the LoRA + RAFE paradigm.
Third, Uni-Mol WL FT Embed results for Flavor remain pending (WL finetuning job not yet
complete); once available they will complement the FL results (MCC~$= 0.765$) already reported
and enable a direct FL-vs-WL comparison for tree classifiers on 3D embeddings in multi-class
taste prediction.

In conclusion, PEARL establishes a practical framework for molecular property
prediction that combines parameter-efficient 3D and SMILES-based finetuning with
retrieval-augmented chemical neighbourhood context.
By achieving state-of-the-art performance on ClinTox toxicity prediction and BACE inhibitor
classification while reducing trainable parameters by 75--97\%, PEARL enables
resource-constrained laboratories to deploy high-performance molecular AI tools.

%% ─────────────────────────────────────────────────────────────────────────────
\section*{Key Points}
%% ─────────────────────────────────────────────────────────────────────────────

\begin{itemize}
  \item PEARL applies LoRA to three CLM families---ChemBERTa, MolFormer, and Uni-Mol---
        achieving 75--97\% parameter reduction while matching or exceeding full finetuning.
  \item Uni-Mol's SE(3)-equivariant 3D representations achieve the best BACE inhibitor
        classification (MCC~$= 0.623$), an 8\% improvement over the best SMILES-based method.
  \item The RAFE framework with ZINC-250k neighbourhood context sets a new state of the art on
        ClinTox toxicity prediction (CT\_TOX MCC~$= 0.913$) and matches E2E LoRA on BBBP
        (MCC~$= 0.798$) with a fully interpretable tree-based pipeline.
  \item MolFormer consistently outperforms ChemBERTa; weighted loss excels on imbalanced
        datasets while focal loss is optimal for balanced, challenging tasks (BACE).
  \item Adding PC features to finetuned embeddings provides minimal benefit, as LoRA-adapted
        models already internalise task-relevant molecular properties.
\end{itemize}

%% ─────────────────────────────────────────────────────────────────────────────
\section*{Acknowledgments}
%% ─────────────────────────────────────────────────────────────────────────────

The authors thank the Qatar National Library (QNL) for their support in providing resources
and funding the APC.

\section*{Funding Sources}
This research did not receive any specific grant from funding agencies in the public,
commercial, or not-for-profit sectors.

\section*{Conflicts of Interest}
The authors declare no competing financial interests. R.M.\ is Scientific Advisor at Eternal,
U.A.E.

\section*{Author Contributions}
K.A.\ and H.H.\ contributed to Methods, Experiments, Validation, Writing -- original draft.
I.S.\ contributed to Investigation, Writing -- review \& editing.
M.C.\ and H.B.\ participated in Investigation, Infrastructure, Writing -- review \& editing.
S.G.\ participated in Conceptualization, Project Supervision, Validation, Writing -- review \& editing.
R.M.\ conceived, conceptualised, supervised the project, contributed in Method, Writing -- review \& editing.

%% ─────────────────────────────────────────────────────────────────────────────
\bibliographystyle{plainnat}
\bibliography{PEARL_references}

%% ─────────────────────────────────────────────────────────────────────────────
\appendix
\section*{Supplementary Material}
%% ─────────────────────────────────────────────────────────────────────────────

\noindent\textbf{Supplementary Table S1.} RDKit-derived molecular descriptors, NetworkX graph
features, Morgan fingerprints, and MACCS keys comprising the 473-dimensional engineered
feature vector.

\noindent\textbf{Supplementary Table S2.} Full BBBP ablation: all modality $\times$ CLM $\times$
loss $\times$ ML combinations.

\noindent\textbf{Supplementary Table S3.} Full Flavor (FART) ablation including per-class AUPR.

\noindent\textbf{Supplementary Table S4.} Full ClinTox ablation for CT\_TOX and FDA\_APPROVED.

\noindent\textbf{Supplementary Table S5.} Full BACE ablation.

\noindent\textbf{Supplementary Table S6.} Complete LoRA parameter efficiency analysis for all
model families, datasets, and loss functions, including Uni-Mol configurations.

\noindent\textbf{Supplementary Section S1.} Zero-shot vs.\ full finetuning: background and
comparison.

\noindent\textbf{Supplementary Section S2.} Detailed LoRA mathematical formulation and
implementation for Uni-Mol (\texttt{in\_proj} targeting, gradient flow preservation,
conformer caching strategy).

\noindent\textbf{Supplementary Section S3.} PEARL custom trainer: loss function
implementations, class-weight computation, and evaluation metric computation.

\noindent\textbf{Supplementary Section S4.} RAFE pipeline details: ZINC-250k cleaning
procedure, FAISS index construction (FlatIP vs.\ HNSW trade-offs), PCA fitting strategy,
cross-space consistency feature computation.

\end{document}
